% H38: platform stalls and scheduler synchrony. Sources: hypotheses/H38-platform-stalls/{README.md (Round 1b), summary/content.tex, summary/meta.json};
% writeup/visuals/H38-platform-stalls/{caption.tex, README.md}; interpretation/swarm-constants.json (f_sched).
\section{H38: The runner, not the platform, makes most synchrony}

\textbf{The question:}
H02 and H19 found collective co-activation: agents are active in the same minutes more often than chance predicts.
Is part of it infrastructure: platform stalls that silence everyone, or a runner that starts and stops all agents together?

\textbf{The intuition:}
Treat each agent as a spin: active or quiet in each minute.
In the Curie--Weiss model each spin couples equally to the total activity $M$, the mean-field coupling.
A stall is a second field, a common input that switches all agents off at once.
A shared on/off field inflates the variance of $M$ even with no coupling.
So the gain $g=1-1/\mathrm{VR}$, built from the variance ratio VR, looks like coupling.
Stalls alone can give $g\approx0.34$ for 12 agents if they fill 5\% of minutes.
The fix is to remove the field before reading the coupling.

\textbf{What we measured:}
H38 builds an outage table: minutes with at most one active agent, with causes from the logs.
The observable is the equal-time co-activation gain in 30-min blocks.
The adjusted gain drops off-schedule minutes and fills in agents that have not started, have finished or are consolidating.
The DQ8 variant trims each day to the window in which all agents run.
The scheduler share is $f=1-E_{\rm adj}/E_{\rm raw}$, where $E$ is the excess gain over the null.
The null shifts each agent's series within 30-min blocks, \emph{after} the trim; this null's false-positive rate is 2--4\%, against 28--34\% on whole-day grids.
Round 1b covers 35 goal periods outside the reserved data on the corrected activity table.
\emph{Scheduler field:} it is the object of the test, and the trim plus conditioning removes it.
\emph{External drive:} partly removed; a 30-min block field absorbs slow drives, but NE14 coincides with a goal change.
\emph{Shared model prior:} open; silences that follow one model provider are not tested.
\emph{Common input:} not relevant, because the card makes no influence claim.
\emph{Synthetic check:} 50 replicates $\times$ 24 conditions.
Planted stalls alone make the raw gain significant in 98--100\% of replicates.
The pre-registered stall filter still reports false coupling in 20--60\% of them.
Agent-state conditioning cuts that to 6--36\% and recovers planted coupling within $-5\%$ to $+23\%$.

\begin{figure}[h]
\includegraphics[width=\textwidth]{H38-platform-stalls/fig.pdf}
\caption{Scheduler or coupling?
In a simulated day, agents that start and stop together inflate the co-activation gain. The gain falls once the day is trimmed to the window in which all agents run.
In the village, most regime-III excess (orange) goes away with this correction, while regime-I excess (blue) stays.
Animation: \texttt{writeup/visuals/H38-platform-stalls/anim.mp4}.}
\label{fig:int-H38}
\end{figure}

\textbf{What we found:}
Joint silences are rare and scheduled.
They fill 5.3\% of minutes; round 1 reported 14.3\%, because the old activity table dropped half of all events.
Of these minutes, 78\% fall outside the operator's run, and infrastructure errors explain 0.1\%.
An error burst makes a joint silence \emph{less} likely (median log odds ratio $-0.96$).
In regime III the scheduler share is $f=0.72$ (8 periods), and 0.83 with the trim.
After the trim only 3/8 regime-III periods keep a significant excess: \#37 ($z=2.0$), \#44 ($z=3.4$) and \#51 ($z=8.5$).
With conditioning on top, only \#44 and \#51 do.
Regime I keeps its excess in 16/18 periods ($f=0.11$).
At NE43 the day-edge share does not drop when the operator's pause and resume messages stop (0.19 to 0.14, $p=0.08$).
So the runner's schedule, not the announcement, makes the edges.
H50 finds the same field with an independent method on call windows.
The claim ``regime-III co-activation is 70--80\% the runner's start/stop schedule'' stands\cred{0.90}.
Its expected usefulness is EU~2.93 (value if true 3.25; two raters, one blind).
Small conflict: \texttt{meta.json} rounds the share to 70--80\%; the card's round 1b gives 0.72 (conditioning) and 0.83 (trim).

\textbf{What it means:}
Advice for an always-on swarm with a daily start and stop: trim each day to the window in which all agents run. Then draw the null samples.
In this village that one step removes about 70--80\% of regime-III co-activation.
Find the edges from the agents' first and last actions, not from operator messages.
Do not alert on joint silences as a stall signal: error logs do not predict them.
Message-driven swarms (regime I) gain nothing from this rule.

\textbf{Caveats:}
A failed LLM call leaves no log line, so provider outages are untested.
On the corrected table most periods have only a few joint-silence minutes.
The headline estimator came after the synthetic smoke runs; the pre-registered stall filter failed.
The \texttt{scheduled} flag comes from operator messages and is empty after 08-05 (inside \#51).
A stall that fills a whole 30-min block leaves no trace in the block-detrended gain.
The fragile flag is set: the coordinator rated the claim fragile.
The confirmation script must be re-frozen on the corrected table, with the trim variants, before it runs.
No confirmation test on reserved goal periods has run.
